% !TEX root = ../nc-quadric-classification.tex
\section{cubic 型による被覆の限界}
\label{sec:cubic-coverage}

\begin{theorem}[cubic 表示の判定]
\label{thm:cubic-coverage}
$\HH$ 対象 $(\cC,\OO)$ の任意の適合表示から
定理~\ref{thm:normal-helix} の二つの初期 MCM 型を取る。
次の条件は同値である。
\begin{enumerate}
\item $(\cC,\OO)$ はある Van den Bergh の三次三次元 AS 正則
$\mathbb Z$-algebra $D$ に対する $(\qgr D,\pi e_0D)$ と pointed exact 同値である。
\item 二つの初期型の少なくとも一方で、偶数添字間の隣接合成が階数四になる。
\end{enumerate}
特に命題~\ref{prop:nonstandard-rank-three-example} の対象は
いかなる cubic 表示でも添字零の構造対象を保って実現できない。
\end{theorem}

\begin{proof}
十分性は定理~\ref{thm:normal-helix} と系~\ref{cor:normal-cubic-defect}
から従う。必要性のため、cubic 表示があると仮定し、
$\pi e_iD$ の像を $Q_i\in\cC$ と書く。
cubic 正則性と \cite[Proposition 4.11, Theorem 6.8]{MN2025} により、
$tQ_i\simeq Q_{i-4}$、および
$\Hom(Q_j,Q_i)=D_{ij}$ である。
この適用に必要な正則性と周期性は \cite[arXiv v3, Definition 4.1.1,
Proposition 5.6.1]{VdB2011} と
\cite[Theorem 4.19]{MN2025} から従う。

$M=Q_{-2}$ と置く。二添字移動は $D$ の自己同型として実現されるので
$Q_{-2}$ は $\OO=Q_0$ の exact 自己同値による像である。
従って $M$ は rank one torsion-free exceptional である。
さらに cubic 型の成分次元から
$\Hom(\OO,t^nM)=D_{-4n-2,0}$ は $(2n+2)^2$ 次元である。命題~\ref{prop:first-orbit-object} の証明の Ext 消滅により、
$n\gg0$ では Euler 特性も同じ値を取る。
命題~\ref{prop:first-orbit-object} により $M\simeq s\OO$ となる。
従ってすべての整数 $j$ で $Q_{-2j}\simeq s^j\OO$ であり、
これらの同型は Hom と合成で定めた二つの添字付き環の同型
$D_{\mathrm{even}}\simeq\check B$ を与える（偶数添字は二で割る）。
ここでは cubic 側の二添字移動そのものが反 Serre の平方根であることは使わない。

$P_i=e_iD$ の偶数成分への制限を $P_i^{\mathrm{ev}}$ と書く。
上の同型により次数付き右 $B$-加群とみなせる。
奇数 $i$ に対する右単純加群の極小分解を偶数成分に制限すると
\[
 0\longrightarrow P_{i+4}^{\mathrm{ev}}
 \longrightarrow (P_{i+3}^{\mathrm{ev}})^2
 \longrightarrow (P_{i+1}^{\mathrm{ev}})^2
 \longrightarrow P_i^{\mathrm{ev}}\longrightarrow0
\]
となる。偶数の添字の制限は $B$ の自由シフトである。
$i$ を前後に四ずつ変えると、$P_i^{\mathrm{ev}}$ は自由加群への埋込みを
繰り返し持ち、その余核は次の自由加群へ埋め込める。
従ってこの列の二回の次元移動を反復し、
$B$ の有限単射次元を使えば $\Ext_B^j(P_i^{\mathrm{ev}},B)=0$ ($j>0$) となる。

$X=P_{-1}^{\mathrm{ev}}$ は次数零に二元生成され、
$H_X(z)=2/(1-z)^3$ であり、自由加群ではない。
定理~\ref{thm:normal-helix} の二型の記述により、非自由 MCM の
直和因子は少なくとも二元を必要とする。従って $X$ は
定理~\ref{thm:normal-helix} の二型の一方である。
$D_{-2,-1}$ は左乗法と制限により $X\to B(1)$ の射を与え、
$D_{-1,0}$ は $B\to X$ の射を与える。
それらの積は $D_{-2,0}\simeq B_1$ を張る。
実際、cubic 型では距離二に関係がなく、この隣接積は同型である。
これらは tails への移行後も同じ Hom と合成になるので、
$X$ の偶数項間の合成は階数四である。
\end{proof}
